% Source: https://artofproblemsolving.com/community/c3h1529352_fun_math_problem
% Stable identity: aops-topic-1529352-q01
\textbf{The limit exists.} For $a=0$ every term is $0$, so $f(0)=0$, which is not positive. Let $a>0$ and let $L=\dfrac{1+\sqrt{1+4a}}{2}$, the positive root of $L^2=a+L$. Then $x_1=\sqrt a<L$, and if $x_k<L$ then $x_{k+1}=\sqrt{a+x_k}<\sqrt{a+L}=L$. Also $x_2=\sqrt{a+\sqrt a}>\sqrt a=x_1$, and if $x_k>x_{k-1}$ then $x_{k+1}=\sqrt{a+x_k}>\sqrt{a+x_{k-1}}=x_k$. So the sequence is increasing and bounded above by $L$, hence it converges to some limit $f$ with $0<f\le L$ and $f=\sqrt{a+f}$, i.e.\ $f^2-f-a=0$. The only positive root is $L$, so
\[
f(a)=\dfrac{1+\sqrt{1+4a}}{2}.
\]

\textbf{Integer values.} $f(a)=n$ for a positive integer $n$ exactly when $n^2-n=a$, i.e.\ $a=n(n-1)$; since $a>0$ we need $n\ge2$. The smallest such $a$ are
\[
2\cdot1=2,\quad3\cdot2=6,\quad4\cdot3=12,\quad5\cdot4=20,\quad6\cdot5=30,
\]
with $f(a)=2,3,4,5,6$. Their sum is $2+6+12+20+30=\boxed{70}$.
